% Part: proof-theory
% Chapter: proof-search
% Section: completeness

\documentclass[../../../include/open-logic-section]{subfiles}

\begin{document}

\olfileid[tr]{pt}{ps}{cpl}

\olsection{Tamlık}

Şimdi ispat arama algoritmasının şaşmaz olduğunu gösteriyoruz; yani algoritma
başarısız olarak durur ya da hiç durmazsa başlangıç sekantı~$\Gamma \Sequent
\Delta$ için !!{proof} yoktur. Bunu yapmak üzere bütün $!A \in \Gamma$
öğeleri için $\Sat{M}{!A}$ ve bütün~$!A \in \Delta$ öğeleri için
$\Sat/{M}{!A}$ olacak biçimde !!a{structure}~$\Struct M$ tanımlarız. Başka
bir deyişle~$\Gamma$ içindeki bütün !!{formula} biçimleri doğru,~$\Delta$ içindeki
bütün !!{formula} biçimleri ise~$\Struct M$ içinde yanlıştır; dolayısıyla $\Gamma
\Sequent \Delta$ geçerli değildir. \Log{G3c} sağlam olduğundan $\Gamma
\Sequent \Delta$ için !!{proof} yoktur.

$\Struct{M}$'yi, !!{formula} biçimlerinden oluşan iki küme~$\Theta$, $\Xi$ ile
!!{constant} simgelerinden oluşan bir küme~$C$ üzerinden tanımlarız. $\Theta$, $\Xi$ ve~$C$'yi,
başarısız olmuş (başarısız olarak durmuş ya da hiç sona ermemiş) ispat
aramasının bir \emph{başarısızlık dalı}ndan elde ederiz. Başarısızlık dalı,
$\Pi_n \Sequent \Lambda_n$ sekantları ve !!{constant} simgesi kümeleri~$C_n$
dizisidir; burada $\Pi_0 \Sequent \Lambda_0$, son sekant~$\Gamma \Sequent
\Delta$'dır ve $\Pi_{n+1} \Sequent \Lambda_{n+1}$, $n+1$.~aşamada $\Pi_n
\Sequent \Lambda_n$ için üretilen ve algoritmanın !!a{proof} bulamadığı bir
öncüldür. Aksi hâlde algoritma !!a{proof} bulmuş olacağından her~$n$ için
böyle bir öncülün var olması gerektiği açıktır. $C_n$,~$\Pi_n \Sequent
\Lambda_n$ ile ilişkilendirilmiş !!{constant} simgesi kümesi olsun.

Algoritma yalnızca atomik !!{formula} biçimleri içeren bir en üst sekantta başarısız
olarak durursa o sekantta biten dal bir başarısızlık dalıdır. Algoritma hiç
sona ermezse gittikçe daha büyük ağaçlar üretmeyi sürdürür. Bunların sonsuz
bir ağaç büyüttüğünü düşünebilirsiniz (algoritma sonsuz sayıda adım çalışmış
olsaydı elde edeceğiniz ağaç). Bu ağaç sonsuz, fakat sonlu dallanan bir ağaç
olurdu (her düğümün yalnızca bir ya da iki çocuğu---hemen üzerindeki
sekantlar---vardır). K\H{o}nig'in lemması gereği her sonsuz ve sonlu dallanan
ağacın sonsuz bir dalı vardır. Arama algoritmasının sonsuza dek çalıştığı
durumda böyle bir sonsuz dal, bir başarısızlık dalı olurdu.

$\Pi_n \Sequent \Lambda_n$ ile $C_n$ dizisi bir başarısızlık dalıysa
$\Theta = \bigcup_n \Pi_n$ ve $\Xi = \bigcup_n \Lambda_n$ (indisleri ve
!!{formula} öğelerinin yinelenen oluşumlarını kaldırarak), ayrıca $C =
\bigcup_n C_n$ olsun.

Artık~$\Struct{M}$'yi tanımlayabiliriz.
\begin{enumerate}
\item Evren $\Domain{M}$,~$C$ ile $\Gamma \Sequent \Delta$ içinde bulunan
!!{function} simgeleri kullanılarak kurulabilen, fakat hiç !!{variable} içermeyen
terimler kümesidir.
\item $c \in \Domain{M}$ ise $\Assign{c}{M} = c$.
\item $f$,~$\Gamma \Sequent \Delta$ içinde geçen $m$-konumlu bir
!!{function} ve $t_1$, \dots, $t_m \in \Domain{M}$ ise
$\Assign{f}{M}(t_1, \dots, t_m) = f(t_1, \dots, t_m)$.
\item $R$,~$\Gamma \Sequent \Delta$ içinde geçen $m$-konumlu bir
!!{predicate} ise $\Assign{R}{M} = \Setabs{\tuple{t_1, \dots, t_m} \in
\Domain{M}^m}{R(t_1, \dots, t_m) \in \Theta}$.
\end{enumerate}
$\Struct{M}$, evreni bir terimler kümesi olduğu ve !!{constant} ile
!!{function} simgelerinin $\Value{t}{M} = t$ olmasını güvence altına alacak biçimde
yorumlandığı için bir \emph{terim modeli}dir. $\Theta$ içindeki atomik
!!{formula} biçimleri, $\Assign{R}{M}$'yi $\Sat{M}{R(t_1, \dots, t_m)}$ ancak ve
ancak $R(t_1, \dots, t_m) \in \Theta$ ise geçerli olacak biçimde tanımlamak
için kullanılır.

\begin{lem}\ollabel{lem:truth}
Bütün $!A \in \Theta \cup \Xi$ için:
\begin{enumerate}
  \item $!A \in \Theta$ ise $\Sat{M}{!A}$.
  \item $!A \in \Xi$ ise $\Sat/{M}{!A}$.
\end{enumerate}
\end{lem}

\begin{proof}
  $\depth{!A}$ üzerinde tümevarımla.

Önce $!A$'nın atomik olduğunu, yani ya $!A \ident \lfalse$ ya da $!A
\ident R(t_1, \dots, t_m)$ olduğunu varsayın.

$\Pi_n \Sequent \Lambda_n$ bir başarısızlık dalı olduğundan hiçbir $\Pi_n$
$\lfalse$ içeremez (aksi hâlde $\Pi_n \Sequent \Lambda_n$ bir aksiyom
olurdu). $\Struct{M}$ tanımı gereği $\Sat{M}{R(t_1, \dots, t_m)}$ ancak ve
ancak $R(t_1, \dots, t_m) \in \Theta$ ise geçerlidir. Bunlar birlikte,
$!A \in \Theta$ ise $\Sat{M}{!A}$ olduğunu gösterir.

$!A$ atomik ve $!A \in \Pi_n$ ise $!A \in \Pi_{n+1}$; $!A \in
\Lambda_n$ ise $!A \in \Lambda_{n+1}$ olur. (Bu, ispat arama
algoritmasının ardıl sekantları üretme biçiminden çıkar.) Dolayısıyla $!A
\in \Theta \cap \Xi$ ise bazı~$n$ için $!A \in \Pi_n \cap \Lambda_n$
olur. Bu, $\Pi_n \Sequent \Lambda_n$'nin bir aksiyom olduğu anlamına gelir;
oysa başarısızlık dalı hiç aksiyom içeremez. Böylece atomik~$!A$ için kuruluş
gereği $!A \notin \Theta \cap \Xi$ olur; bunun sonucu olarak $!A \in \Xi$
ise $!A \notin \Theta$'dır. $\Struct{M}$ tanımı gereği bu, $!A \in \Xi$
ise $\Sat/{M}{!A}$ olduğunu gösterir.

Tümevarım adımı için (1) ve (2)'nin bütün !!{formula} biçimleri için,
bu biçimler $!A$'dan daha düşük !!{depth} taşıdığında geçerli olduğunu
varsayın. Durumları ayırarak (1) ve
(2)'yi~$!A$ için kanıtlarız.

$!A^i$,~$\Pi_n \Sequent \Lambda_n$ içinde geçtiğinde $i$,~$\Pi_n \Sequent
\Lambda_n$ içindeki en küçük indis olmadığı sürece $!A^i$'nin~$\Pi_{n+1}
\Sequent \Lambda_{n+1}$ içinde de geçtiğine önce dikkat edin. Her aşamada
eklenen en üst sekantlarda en küçük indis büyüdüğünden, sonunda $!A^i$'nin
geçtiği ve $i$'nin en küçük indis olduğu bir $\Pi_n \Sequent \Lambda_n$
sekantına ulaşırız. Algoritma $n+1$.~aşamada~$!A^i$'yi indirger. Başka bir
deyişle $!A \in \Theta$ ise bazı $n$ için $\Pi_n = \Pi_n', !A^i$ ve $i$ en
küçük indistir. $!A \in \Xi$ ise bazı $n$ için $\Lambda_n = \Lambda_n',
!A^i$ ve~$i$ en küçük indistir.

\begin{enumerate}
  \item $!A \in \Theta$ ve $!A \ident !B \land !C$. O zaman bazı $n$ için
  $\Pi_n = \Pi_n', (!B \land !C)^i$. $\Pi_{n+1} \Sequent \Lambda_{n+1}$,
  \LeftR{\land} kuralının karşılık gelen öncülüdür; yani $\Pi_{n+1} =
  \Pi_n', !B^{k+1}, !C^{k+2}$. Dolayısıyla $!B, !C \in \Theta$.
  Tümevarım varsayımı gereği $\Sat{M}{!B}$ ve $\Sat{M}{!C}$. $\Sat{M}{}$
  tanımı gereği $\Sat{M}{!B \land !C}$ elde edilir.

  \item $!A \in \Xi$ ve $!A \ident !B \land !C$. O zaman bazı $n$ için
  $\Lambda_n = \Lambda_n', (!B \land !C)^i$. $\Pi_{n+1} \Sequent
  \Lambda_{n+1}$, sonucu $\Pi_n \Sequent \Lambda'_n, (!B \land !C)^i$ olan
  \RightR{\land} kuralının iki öncülünden biridir; yani $\Lambda_{n+1} =
  \Lambda_n', !B^{k+1}$ ya da $\Lambda_{n+1} = \Lambda_n', !C^{k+1}$.
  Dolayısıyla ya $!B \in \Xi$ ya da $!C \in \Xi$. Tümevarım varsayımı
  gereği ya $\Sat/{M}{!B}$ ya da $\Sat/{M}{!C}$. $\Sat{M}{}$ tanımı gereği
  $\Sat/{M}{!B \land !C}$ elde edilir.

  \item $!A \ident \lnot !B$, $!A \ident !B \lor !C$ ve $!A \ident !B
  \lif !C$ durumları benzerdir ve alıştırma olarak bırakılmıştır.

  \item $!A \in \Theta$ ve $!A \ident \lexists[x][!B(x)]$. O zaman bazı
  $n$ için $\Pi_n = \Pi_n', \lexists[x][!B(x)]^i$. $\Pi_{n+1} \Sequent
  \Lambda_{n+1}$, \LeftR{\lexists} kuralının karşılık gelen öncülüdür;
  yani bazı~$c$ için $\Pi_{n+1} = \Pi_n', !B(c)^{k+1}$. Dolayısıyla
  $!B(c) \in \Theta$ ve~$c \in C$. Tümevarım varsayımı gereği
  $\Sat{M}{!B(c)}$. $\Sat{M}{}$ tanımı gereği
  $\Sat{M}{\lexists[x][!B(x)]}$ elde edilir.

  \item $!A \in \Xi$ ve $!A \ident \lexists[x][!B(x)]$. O zaman bazı $n$
  için $\Lambda_n = \Lambda_n', \lexists[x][!B(x)]^i$. $\lexists[x][!B(x)]$
  her indirgendiğinde yeni bir indisle öncül sekantının sağında kaldığından,
  bir başarısızlık dalında sağda geçiyorsa sağ tarafta sonsuz kez indirgenir.
  Her $t \in \Domain{M}$ sonunda bir başarısızlık dalında ``yalnızca~$C_n$
  içindeki sabitleri içeren ve $\lexists[x][!B(x)]$ sağda indirgenirken daha
  önce kullanılmamış ilk terim'' olur. Dolayısıyla bütün $t \in \Domain{M}$
  için bazı~$n$ bakımından $!B(t) \in \Lambda_n$ ve sonuç olarak $!B(t) \in
  \Xi$ olur. Tümevarım varsayımı gereği bütün $t \in \Domain{M}$ için
  $\Sat/{M}{!B(t)}$. $\Sat{M}{}$ tanımı gereği
  $\Sat/{M}{\lexists[x][!B(x)]}$ elde edilir.

  \item $!A \ident \lforall[x][!B(x)]$ durumları benzerdir ve alıştırma
  olarak bırakılmıştır.
\end{enumerate}
\end{proof}

\begin{cor}\ollabel{cor:complete} \Log{G3c} için ispat arama algoritması
tamdır: Başarısız olarak durur ya da hiç sona ermezse başlangıç
sekantı~$\Gamma \Sequent \Delta$ geçerli değildir ve dolayısıyla !!{proof}
yoktur.
\end{cor}

\begin{proof}
  Algoritma başarısız olarak durur ya da hiç sona ermezse, üzerinden
  $\Struct{M}$ tanımlanabilen bir başarısızlık dalı vardır. \olref{lem:truth}
  gereği bütün $!A \in \Gamma$ için $\Sat{M}{!A}$ ve bütün $!A \in \Delta$
  için $\Sat/{M}{!A}$. ($\Pi_0 = \Gamma$ ve $\Lambda_0 = \Delta$ olduğunu,
  dolayısıyla $\Gamma \subseteq \Theta$ ile $\Delta \subseteq \Xi$ olduğunu
  hatırlayın.) Öyleyse $\Gamma \Sequent \Delta$ geçerli değildir.
  \Log{G3c}'nin sağlamlığı gereği $\Gamma \Sequent \Delta$ için !!{proof}
  yoktur.
\end{proof}

\begin{cor}
$\Gamma \Sequent \Delta$ geçerliyse $\Log{G3c} \Proves \Gamma \Sequent
\Delta$ olduğundan \Log{G3c} tamdır.
\end{cor}

\begin{proof}
  Karşıt tersini kanıtlarız. $\Log{G3c} \Proves/ \Gamma \Sequent \Delta$
  olduğunu varsayın. Bulunacak !!{proof} olmadığından ispat arama algoritması
  başarısız olarak durmalı ya da hiç sona ermemelidir. \olref{cor:complete}
  gereği $\Gamma \Sequent \Delta$ geçerli değildir.
\end{proof}


\end{document}
